\chapter{机械结构与整机安装}

\section{无人机总体结构}

说明机架构型、轴距、材料、桨叶尺寸、电机与保护罩选择，并结合比赛尺寸限制解释选型依据。

\placeholderfigure{uav-overview.pdf}{无人机整机俯视、前视与侧视图}

\todoitem{仓库尚未包含完整机械/CAD资料。补充机架型号、轴距、整机外廓尺寸、质量、重心、惯量估计和有效载荷。}

\section{传感器与飞控板安装}

重点说明 ICM42688、ToF、光流和摄像头相对机体系的位置与朝向，给出安装矩阵、支架尺寸及减振方案。
IMU 安装矩阵已有六面标定流程，可在本节说明机械安装误差如何通过标定矩阵补偿。

\placeholderfigure{sensor-layout.pdf}{传感器、飞控板及坐标轴安装关系}

\section{摄像头布局与视场覆盖}

说明双摄像头的安装高度、俯仰角、航向偏置、重叠视场和盲区，并将机械参数与第五章的相机模型对应。

\todoitem{给出相机到机体系的平移向量与旋转矩阵实测值；用赛场俯视图验证覆盖范围。}

\section{线束、供电与空地连接}

说明动力线、信号线、ToF/光流接口、地面车模通信与可能存在的系留供电路径。大电流线束和敏感信号线应分开布置，
并给出插接件防松、应力释放及紧急断开设计。

\section{地面车模接口}

本仓库主要记录无人机端，地面车模机械结构需与车模项目联合补充，包括车模尺寸、转向/驱动结构、车载标志、
无人机承载或供电结构以及踩灯/盲盒执行机构。

\placeholderfigure{ground-car-overview.pdf}{地面车模与空地接口结构}

\section{机械参数与可靠性验证}

\begin{table}[htbp]
  \centering
  \caption{机械结构待测参数}
  \begin{tabularx}{\textwidth}{p{4.3cm}p{3.2cm}X}
    \toprule
    参数 & 实测值 & 测量方法/说明 \\
    \midrule
    无人机长×宽×高 & 待实测 & 含保护罩与突出部件 \\
    无人机起飞质量 & 待实测 & 含电池、线束和全部传感器 \\
    地面车模长×宽×高 & 待实测 & 按比赛测量口径 \\
    地面车模总质量 & 待实测 & 含任务装置 \\
    重心偏移 & 待实测 & 相对机架几何中心 \\
    一阶结构振动频率 & 待测试 & 结合IMU频谱或模态测试 \\
    \bottomrule
  \end{tabularx}
\end{table}
